% from ega0-1-fr.tex
\oldpage[0\textsubscript{I}]{33}

Cela résulte en effet aussitôt des définitions, tenant compte de ce que
$\psi''_*=\psi'_*\circ\psi_*$. En outre, si
$u:\mathcal{G}\to\psi_*(\mathcal{F})$ est un $\psi$-morphisme,
$v:\mathcal{H}\to\psi'_*(\mathcal{G})$ un $\psi'$-morphisme, et
$w=\psi'_*(u)\circ v$ leur composé (3.5.2), on constate aussitôt que
$w^\sharp$ est le morphisme composé
\[
  w^\sharp:\psi^*(\psi'^*(\mathcal{H}))
  \xrightarrow{\psi^*(v^\sharp)}
  \psi^*(\mathcal{G})
  \xrightarrow{u^\sharp}\mathcal{F}.
\]
\end{env}

\begin{env}[3.5.6]
\label{0.3.5.6-fr}
Prenons en particulier pour $\psi$ l'application identique
$1_X:X\to X$. Alors si l'image réciproque par $\psi$ d'un préfaisceau
$\mathcal{F}$ sur $X$ à valeurs dans $\mathbf{K}$ existe, on dit que cette
image réciproque est le \emph{faisceau associé au préfaisceau $\mathcal{F}$}.
Tout morphisme $u:\mathcal{F}\to\mathcal{F}'$ de $\mathcal{F}$ dans un
\emph{faisceau} $\mathcal{F}'$ à valeurs dans $\mathbf{K}$ se factorise donc
de façon unique en
\[
  \mathcal{F}\xrightarrow{\rho_{\mathcal{F}}}
  1_X^*(\mathcal{F})\xrightarrow{u^\sharp}\mathcal{F}'.
\]
\end{env}

\subsection{Faisceaux simples et faisceaux localement simples}

\begin{env}[3.6.1]
\label{0.3.6.1-fr}
Nous dirons qu'un \emph{préfaisceau} $\mathcal{F}$ sur $X$, à valeurs dans
$\mathbf{K}$, est \emph{constant} si les morphismes canoniques
$\mathcal{F}(X)\to\mathcal{F}(U)$ sont des \emph{isomorphismes} pour tout
ouvert non vide $U\subset X$~; on notera que $\mathcal{F}$ n'est pas
nécessairement un faisceau. On dit qu'un \emph{faisceau} est \emph{simple}
s'il est associé (3.5.6) à un préfaisceau constant. On dit qu'un faisceau
$\mathcal{F}$ est \emph{localement simple} si tout $x\in X$ admet un voisinage
ouvert $U$ tel que $\mathcal{F}|U$ soit simple.
\end{env}

\begin{env}[3.6.2]
\label{0.3.6.2-fr}
Supposons que $X$ soit \emph{irréductible} (2.1.1)~; alors les propriétés
suivantes sont équivalentes~:
\begin{enumerate}
  \item[a)] \emph{$\mathcal{F}$ est un préfaisceau constant sur $X$}~;
  \item[b)] \emph{$\mathcal{F}$ est un faisceau simple sur $X$}~;
  \item[c)] \emph{$\mathcal{F}$ est un faisceau localement simple sur $X$}.
\end{enumerate}

En effet, soit $\mathcal{F}$ un préfaisceau constant sur $X$~; si $U$, $V$
sont deux ouverts non vides dans $X$, $U\cap V$ est non vide, donc
$\mathcal{F}(X)\to\mathcal{F}(U)\to\mathcal{F}(U\cap V)$ et
$\mathcal{F}(X)\to\mathcal{F}(U)$ étant des isomorphismes, il en est de même
de $\mathcal{F}(U)\to\mathcal{F}(U\cap V)$ et de même
$\mathcal{F}(V)\to\mathcal{F}(U\cap V)$ est un isomorphisme. On en conclut
aussitôt que l'axiome (F) de (3.1.2) est bien vérifié, $\mathcal{F}$ est
isomorphe à son faisceau associé, et par suite a) entraîne b).

Soit maintenant $(U_\alpha)$ un recouvrement ouvert de $X$ par des ouverts
non vides et un faisceau $\mathcal{F}$ sur $X$ tel que
$\mathcal{F}|U_\alpha$ soit simple pour tout $\alpha$~; comme $U_\alpha$ est
irréductible, $\mathcal{F}|U_\alpha$ est un préfaisceau constant en vertu de
ce qui précède. Comme $U_\alpha\cap U_\beta$ n'est pas vide,
$\mathcal{F}(U_\alpha)\to\mathcal{F}(U_\alpha\cap U_\beta)$ et
$\mathcal{F}(U_\beta)\to\mathcal{F}(U_\alpha\cap U_\beta)$ sont des
isomorphismes, d'où on déduit un isomorphisme canonique
$\theta_{\alpha\beta}:\mathcal{F}(U_\alpha)\to\mathcal{F}(U_\beta)$ pour tout
couple d'indices. Mais alors, si on applique la condition (F) pour $U=X$, on
voit que pour tout indice $\alpha_0$, $\mathcal{F}(U_{\alpha_0})$ et les
$\theta_{\alpha_0\alpha}$ sont solutions du problème universel, ce qui (en
vertu de l'unicité) implique que
$\mathcal{F}(X)\to\mathcal{F}(U_{\alpha_0})$ est un isomorphisme, et démontre
donc que c) entraîne a).
\end{env}

\subsection{Images réciproques de préfaisceaux de groupes ou d'anneaux}

